\documentclass[11pt]{article}
\usepackage[margin=1in]{geometry}
\usepackage{amsmath,amssymb,amsthm,hyperref,microtype}
\hypersetup{colorlinks=true,linkcolor=blue,citecolor=blue,urlcolor=blue,pdftitle={Sharp disagreement between egalitarian core selections in assignment games},pdfauthor={Maxwell Flanagan}}
\newtheorem{theorem}{Theorem}
\newtheorem{lemma}[theorem]{Lemma}
\newtheorem{proposition}[theorem]{Proposition}
\newcommand{\Lmin}{\operatorname{Lmin}}
\newcommand{\Lmax}{\operatorname{Lmax}}
\title{Sharp disagreement between egalitarian\\ core selections in assignment games}
\author{Maxwell Flanagan}
\date{September 4, 2026 --- Internal research draft}
\begin{document}
\maketitle
\begin{abstract}
For weighted bipartite assignment games on $n\ge4$ players, we determine
the maximum normalized $\ell_1$ distance between the full leximin and
leximax-min core allocations:
\[
 s_n^{\rm assign}=\frac{\lfloor n/2\rfloor}{\lfloor n/2\rfloor+1}.
\]
The main structural lemma states that leximin gives each player at least
half their payoff in every other core allocation. Matching complementarity
and a minimax directional witness yield the sharp finite bound. A complete
market with identical rows and weights in $\{1,2\}$ attains it. Exact
assignment games instead have coincident selections, a consequence of
known pairwise egalitarian structure. External correctness and novelty
review remain pending.
\end{abstract}

\section{Definitions and main result}
Let $G=(L,R,E)$ be a finite bipartite graph with nonnegative real edge
weights $w$. Its assignment game values a coalition by the maximum matching
weight in its induced subgraph. Write $n=|L|+|R|$ and let $T>0$ be the
maximum matching weight in $G$. Its entire core is
\begin{equation}\label{eq:core}
 C=\left\{p\ge0:\ \sum_a p_a=T,\quad p_u+p_v\ge w_{uv}
          \ \text{for every }uv\in E\right\}.
\end{equation}
The edge inequalities imply all coalition inequalities by summing over a
coalition's optimal matching; the reverse follows from two-player
coalitions. Nonemptiness follows from bipartite matching LP duality.

Let $x=\Lmin(C)$ lexicographically maximize the increasingly sorted
payoffs, and let $z=\Lmax(C)$ lexicographically minimize the decreasingly
sorted payoffs. Retaining players whose core payoff is identically zero
does not change the other coordinates of either full selection. Define
\[
 \rho(C)=\frac{\|x-z\|_1}{T},\qquad
 s_n^{\rm assign}=\sup_{G:\,|L|+|R|=n,\,T>0}\rho(C).
\]
Isolated players are allowed, and $n$ counts them. No exact-game hypothesis
is imposed.
\begin{theorem}\label{thm:sharp}
Fix any optimal matching, delete its zero-weight edges, and let $k\ge1$
be its number of remaining edges. For full leximin $x$ and \emph{any}
core allocation $z$ minimizing the largest payoff,
\[
 \frac{\|x-z\|_1}{T}\le\frac{k}{k+1}.
\]
For $k=1$ the selections coincide. For every $k\ge2$ the bound is attained
by a complete $k\times k$ market with weights in $\{1,2\}$. Hence
$s_n^{\rm assign}=0$ for $n=2,3$, and
$s_n^{\rm assign}=\lfloor n/2\rfloor/(\lfloor n/2\rfloor+1)$ for $n\ge4$.
In particular, $s_n^{\rm assign}\to1$.
\end{theorem}
The result applies to nonunique optimal matchings and minimax faces.
Full leximin is essential; a first-stage maximin point need not suffice.

\section{A simultaneous half-payoff guarantee}
Fix the matching from Theorem~\ref{thm:sharp}, with positive edge weights
$s_1,\ldots,s_k$. Summing its inequalities in \eqref{eq:core} and using
efficiency shows that every matched pair has total $s_i$, and every
unmatched player receives zero throughout the core. These identities
require no nondegeneracy or exactness assumptions.
\begin{lemma}\label{lem:half}
For every $y\in C$ and every player $a$, $y_a\le2x_a$.
Consequently $x_a\ge\tfrac12\max_{y\in C}y_a$ simultaneously for all
players, and $\|x-y\|_1/T\le1$ for every $y\in C$.
\end{lemma}
\begin{proof}
Put $d=y-x$. Suppose first that $a\in L$ and $d_a=\delta>0$. Set
\[
 A=\{u\in L:d_u\ge\delta\},\qquad
 B=\{v\in R:d_v\le-\delta\}.
\]
Matched-pair totals give opposite changes across each matching edge;
all unmatched changes vanish. Thus the matching bijects $A$ and $B$.
At an edge $uv$ tight at $x$, feasibility of $y$ implies
$d_u+d_v\ge0$. Such an edge cannot have $u\notin A$ and $v\in B$.
Therefore the direction $h=\mathbf1_A-\mathbf1_B$ preserves the total
and harms no tight edge inequality. Every decreasing coordinate satisfies
$x_v\ge\delta$ by nonnegativity of $y_v$. Finitely many slack constraints
and positive decreasing coordinates ensure $x+\epsilon h\in C$ for all
sufficiently small $\epsilon>0$.

If $\min_Bx>\min_Ax$, a sufficiently small such step raises the lowest
affected coordinates and decreases only strictly larger coordinates,
strictly improving the ascending sorted vector. This contradicts full
leximin optimality, including when smaller fixed coordinates occur
elsewhere. Hence
\[
 \delta\le\min_Bx\le\min_Ax\le x_a,
 \qquad y_a=x_a+\delta\le2x_a.
\]
Interchange the two market sides for positive right-side changes.
Nonpositive changes satisfy the claim immediately. Finally
$\min(x_a,y_a)\ge y_a/2$, so the overlap is at least $T/2$ and the
normalized distance is at most one.
\end{proof}
The constant two in the coordinatewise comparison is tight: on a
single positive edge, leximin splits its value equally while a core
endpoint gives its entire value to one player.

\section{The sharp finite bound and its attainment}
\begin{proof}[Proof of Theorem~\ref{thm:sharp}: upper bound]
On matched pair $i$, write the change from $x$ to $z$ as $(d_i,-d_i)$,
and put
\[
 q_i=|d_i|,\quad D=\sum_iq_i,\quad M=\max_a z_a>0.
\]
Thus $\|x-z\|_1=2D$. Both $x$ payoffs in pair $i$ are at least $q_i$:
the decreasing endpoint by nonnegativity of $z$, and the increasing
endpoint by Lemma~\ref{lem:half}. The increasing endpoint's $z$ payoff
is at least $2q_i$. Therefore
\begin{equation}\label{eq:pair}
 s_i\ge2q_i,\qquad q_i\le M/2,\qquad D\le kM/2.
\end{equation}
These statements also hold when $q_i=0$.

Some player $a$ satisfies $z_a=M$ and $x_a\ge M$. Otherwise, a
sufficiently short move from $z$ toward $x$ lowers every maximum
coordinate and leaves every other coordinate below $M$, contradicting
minimax optimality. Since $M>0$, this player belongs to a positive
matched pair $p$. Its displacement is $q_p=x_a-M$; the other $x$
payoff in that pair is at least $q_p$. Thus $s_p\ge M+2q_p$.
Summing this strengthened inequality with \eqref{eq:pair} for the other
pairs gives
\[
 T\ge2D+M\ge2D\left(1+\frac1k\right).
\]
This proves the bound, also when $D=0$.

If $k=1$, all other payoffs are fixed at zero. The remaining two sum
to $T$ and range over a closed interval of splits. Maximizing their
minimum is the same as minimizing their maximum, and uniquely selects
the feasible split closest to equality. Hence $x=z$ in this case.
\end{proof}
\begin{proof}[Attainment and player count]
For $k\ge2$, take a complete $k\times k$ market with identical rows
\[
 (w_{i1},\ldots,w_{ik})=(1,\ldots,1,2).
\]
Every perfect matching has value $T=k+1$, and every edge belongs to one.
Complementarity consequently makes every edge inequality tight throughout
the core. Writing $c=(1,\ldots,1,2)$, the entire core is exactly
\[
 u_i=t\quad(1\le i\le k),\qquad
 v_j=c_j-t\quad(1\le j\le k),\qquad0\le t\le1.
\]
Its minimum coordinate is $\min(t,1-t)$, uniquely maximized at $t=1/2$.
Its maximum is $2-t$, uniquely minimized at $t=1$. Every coordinate
changes by $1/2$, so $\|x-z\|_1=k$ and $\rho=k/(k+1)$.
Adding one isolated player leaves the full selectors otherwise unchanged.
Every matching on $n$ players has at most $\lfloor n/2\rfloor$ edges,
which proves the exact $n$-player formula.
\end{proof}

\section{Why exact assignment games have zero divergence}
A game is exact when each coalition's worth is attained by some core
allocation. The following known structural implication explains the
boundary with unrestricted exact TU games; it is not claimed as new
\cite{llerena,solymosi-raghavan}.
\begin{proposition}\label{prop:exact}
If the matched-pair halfsplit $h$ belongs to an assignment-game core,
then $\Lmin(C)=\Lmax(C)=h$. Every exact assignment game satisfies this
hypothesis.
\end{proposition}
\begin{proof}
For any core allocation $y$, let $P$ swap the endpoints of every fixed
matched pair and fix unmatched coordinates. Then $h=(y+Py)/2$.
Thus $h$ is majorized by $y$: pairwise averaging can only increase sums
of the smallest coordinates and decrease sums of the largest coordinates.
If $y\ne h$, strict convexity of squared norm implies
$\|h\|_2^2<\|y\|_2^2$, so their sorted vectors cannot coincide.
Therefore $h$ strictly improves both relevant lexicographic orders.

For an exact assignment game, the two side coalitions $L$ and $R$ have
worth zero and each attains that worth at a core allocation.
Nonnegativity and matched-pair totals force these two allocations to
put every $s_i$ entirely on opposite sides. Their midpoint is $h\in C$.
\end{proof}
The attaining markets above are nonexact: their distinguished right
player always receives at least one, despite having singleton worth zero.
Thus the comparison is between two different restrictions of cooperative
games: general assignment games have maximum divergence tending to one;
within exact assignment games it is zero.

\section{Evidence, related work, and review status}
The argument above is a conventional all-parameter proof. Exact rational
replays reconstruct all coalition values by enumerating matchings,
verify all full selector stages for 255 nonzero $2\times2$ matrices
with entries $0,1,2,3$, and check the half-payoff guarantee by separately
maximizing each coordinate over the entire core. They also exclude any
violating minimax allocation, test rectangular and degenerate examples,
and reconstruct attaining segment cores through $k=6$.
A separate implementation verifies ten even/odd attaining-family cases.
Negative controls distinguish full leximin from first-stage maximin and
minimax comparison points from arbitrary core allocations. The code and
exact outputs accompany this note; finite tests do not replace its proof.

Sorted egalitarian selections and alternating-adjustment methods have
older assignment-game precedents \cite{klijn}; recent work develops
fair-core algorithms \cite{vazirani}. The fair-price midpoint already
provides a half-core-maximum benchmark for another solution concept
\cite{solymosi}. Our candidate contribution is the guarantee specifically
for full leximin and the sharp finite distance to minimax selections.
A dated primary-source audit located no equivalent quantitative theorem,
but does not certify novelty. The audit also records an apparent mismatch
between the graph and printed vectors of Example~1 in Vazirani's v6,
with exact reconstructed selectors and an independent public implementation
fixture corroborating the reconstructed leximin. This discrepancy is
separate from the proof above. External correctness, prior-art, and
significance review remain pending.

\begin{thebibliography}{9}
\bibitem{vazirani} V. V. Vazirani.
\emph{Fair Core Imputations for the Assignment Game: New Solution Concepts
and Efficient Algorithms}. arXiv:2402.11437v6 (2025).
\url{https://arxiv.org/abs/2402.11437}.
\bibitem{klijn} F. Klijn et al.
\emph{Neighbour Games and the Leximax Solution}.
Mathematical Methods of Operations Research 58 (2003), 191--208.
\url{https://pure.uvt.nl/ws/files/533861/110.pdf}.
\bibitem{llerena} F. Llerena.
\emph{The pairwise egalitarian solution for the assignment game}.
Operations Research Letters 40 (2012).
\url{https://doi.org/10.1016/j.orl.2011.11.011}.
\bibitem{solymosi-raghavan} T. Solymosi and T. E. S. Raghavan.
\emph{Assignment games with stable core}.
International Journal of Game Theory 30 (2001), 177--185.
\url{https://doi.org/10.1007/s001820100072}.
\bibitem{solymosi} T. Solymosi.
\emph{Sensitivity of fair prices in assignment markets}.
Mathematical Social Sciences 126 (2023), 1--12.
\url{https://doi.org/10.1016/j.mathsocsci.2023.09.003}.
\end{thebibliography}
\end{document}
